\documentclass[12pt]{article}
\usepackage{amsmath, amssymb, amsthm}
\usepackage{geometry}
\usepackage{hyperref}
\usepackage{listings}
\usepackage{courier}
\usepackage{setspace}
\usepackage{titlesec}

\geometry{margin=1in}
\setstretch{1.15}

\titleformat{\section}{\large\bfseries}{\thesection}{0.5em}{}

\lstset{
    basicstyle=\footnotesize\ttfamily,
    breaklines=true,
    frame=single,
    columns=fullflexible
}

\title{Exploratory Thoughts: Unified Recommender System Architecture\\
\large Transformer Sequence Modeling + Feature Interaction\\
\normalsize \href{https://algo.qq.com/\#intro}{algo.qq.com}}
\author{David \and Xia}
\date{April 2026 \\ \small Draft v0.1}

\begin{document}
\maketitle

\begin{center}
    \small\textit{Prepared in the context of the QQ Algorithm Challenge ---
    \href{https://algo.qq.com/\#intro}{algo.qq.com}}
\end{center}

\begin{abstract}
These are first-thoughts towards a unified recommender system architecture
that merges transformer-based sequence modelling with structured feature
interaction, targeting the QQ Algorithm Challenge. The intent is exploratory:
to sketch the mathematical and implementation foundations before committing to
a full proposal.
\end{abstract}

\noindent\textbf{Keywords:} recommender systems, transformers, feature interaction, PyTorch, deep learning

\bigskip

\section{Overview}
This document presents a compact technical model describing how a unified
recommender-system architecture can merge:
\begin{itemize}
    \item Transformer-style sequence modeling,
    \item Feature-interaction modeling (e.g.\ DeepFM, DCN, DIN),
    \item A unified PyTorch implementation,
    \item A training pipeline and framework comparison.
\end{itemize}

The goal is to show how modern recommender systems integrate temporal
dependencies and structured feature interactions into a single scalable model.

\section{1.\ Minimal Unified Architecture (PyTorch)}
Below is a minimal PyTorch model that combines:
\begin{itemize}
    \item a Transformer encoder for sequential user behaviour,
    \item static feature embeddings (user, item, context),
    \item an MLP for feature interactions.
\end{itemize}

\subsection*{Model Code}
\begin{lstlisting}[language=Python]
import torch
import torch.nn as nn
import torch.nn.functional as F

class UnifiedRecModel(nn.Module):
    def __init__(
        self,
        num_items,
        num_users,
        num_context_feats,
        d_model=64,
        n_heads=4,
        d_ff=128,
        seq_max_len=50,
        mlp_hidden=128,
    ):
        super().__init__()

        self.item_emb = nn.Embedding(num_items, d_model)
        self.user_emb = nn.Embedding(num_users, d_model)
        self.ctx_emb = nn.Embedding(num_context_feats, d_model)

        encoder_layer = nn.TransformerEncoderLayer(
            d_model=d_model,
            nhead=n_heads,
            dim_feedforward=d_ff,
            batch_first=True,
        )
        self.seq_encoder = nn.TransformerEncoder(encoder_layer, num_layers=2)

        in_dim = d_model * 4
        self.mlp = nn.Sequential(
            nn.Linear(in_dim, mlp_hidden),
            nn.ReLU(),
            nn.Linear(mlp_hidden, 1),
        )

    def forward(self, seq_items, user_ids, target_items, ctx_feats, seq_mask=None):
        seq_emb = self.item_emb(seq_items)

        if seq_mask is not None:
            src_key_padding_mask = (seq_mask == 0)
        else:
            src_key_padding_mask = None

        seq_encoded = self.seq_encoder(
            seq_emb,
            src_key_padding_mask=src_key_padding_mask
        )

        if seq_mask is not None:
            lengths = seq_mask.sum(dim=1, keepdim=True).clamp(min=1)
            seq_pooled = (seq_encoded * seq_mask.unsqueeze(-1)).sum(dim=1) / lengths
        else:
            seq_pooled = seq_encoded.mean(dim=1)

        user_vec = self.user_emb(user_ids)
        item_vec = self.item_emb(target_items)
        ctx_vec = self.ctx_emb(ctx_feats).mean(dim=1)

        x = torch.cat([seq_pooled, user_vec, item_vec, ctx_vec], dim=-1)
        logit = self.mlp(x).squeeze(-1)
        return logit
\end{lstlisting}

\section{2.\ Merging Attention and Feature-Interaction Layers}
A unified architecture combines:
\begin{itemize}
    \item Transformer attention for temporal modelling,
    \item Feature-interaction layers (e.g.\ FM, DeepFM) for structured fields.
\end{itemize}

\subsection*{Attention}
\[
\mathrm{Attn}(Q,K,V)
= \mathrm{softmax}\!\left(\frac{QK^\top}{\sqrt{d_k}}\right)V.
\]

\subsection*{Factorization Machine Interaction}
\[
y_{\mathrm{FM}}
= \frac{1}{2}\sum_{i<j}
\left(
\langle v_i, v_j\rangle
\right).
\]

\subsection*{FM Layer Code}
\begin{lstlisting}[language=Python]
class FMInteraction(nn.Module):
    def __init__(self, dim):
        super().__init__()
        self.dim = dim

    def forward(self, field_embs):
        sum_emb = field_embs.sum(dim=1)
        sum_sq = sum_emb * sum_emb
        sq_sum = (field_embs * field_embs).sum(dim=1)
        fm = 0.5 * (sum_sq - sq_sum).sum(dim=1)
        return fm
\end{lstlisting}

\section{3.\ Training Pipeline}
A typical training pipeline includes:
\begin{enumerate}
    \item Data preparation (sequence padding, masks, feature IDs).
    \item PyTorch \texttt{Dataset} and \texttt{DataLoader}.
    \item Training loop with BCE loss.
    \item Validation using AUC and log-loss.
\end{enumerate}

\subsection*{Training Loop}
\begin{lstlisting}[language=Python]
optimizer = torch.optim.Adam(model.parameters(), lr=1e-3)
criterion = nn.BCEWithLogitsLoss()

for epoch in range(num_epochs):
    model.train()
    for batch in train_loader:
        logits = model(
            batch["seq_items"].to(device),
            batch["user_ids"].to(device),
            batch["target_items"].to(device),
            batch["ctx_feats"].to(device),
            batch["seq_mask"].to(device),
        )
        loss = criterion(logits, batch["labels"].float().to(device))
        optimizer.zero_grad()
        loss.backward()
        optimizer.step()
\end{lstlisting}

\section{4.\ PyTorch vs.\ TensorFlow}
\begin{itemize}
    \item \textbf{PyTorch}: flexible, research-friendly, strong Transformer ecosystem.
    \item \textbf{TensorFlow}: strong for production CTR/CVR models, TF-Serving.
    \item \textbf{Deployment}: ONNX, TorchScript, TensorRT, TF-Serving.
    \item \textbf{Trend}: PyTorch dominates research; TF remains strong in industry.
\end{itemize}

\end{document}
